% !TeX root = ../main.tex
\chapter*{Along the Value Stream: What Will Apps Look Like When LLM-OS Arrives?}
\markboth{Along the Value Stream: What Will Apps Look Like When LLM-OS Arrives?}{}
\phantomsection
\addcontentsline{toc}{chapter}{Along the Value Stream: What Will Apps Look Like When LLM-OS Arrives?}
\label{ch:llm_os}

{\kaiti
LLM-OS (Large Language Model Operating System)\footnote{XunJiang. \href{https://xun140202.substack.com/p/along-the-value-stream-what-will}{\textit{Along the Value Stream: What Will App Look Like When LLMOS Arrives?}}, Substack, July 28, 2025.} is on the horizon. Do ordinary developers still have a chance?

If so, what kind of apps can survive?

Written for developers and product people who are still searching for direction but refuse to be marginalized by LLM-OS.
}

\vspace{1em}

Three questions before us:

1. Can a super app still emerge outside of LLM-OS? If so, where does the opportunity lie?

2. If developing for agents is the foreseeable future, should we invest directly in open-source infrastructure, or should we quietly incubate within vertical domains?

3. Choice matters more than effort---what kind of scenario should you choose as your moat?

\section{Why Think About These Questions?}

The people who can join the model teams at ByteDance, DeepSeek, OpenAI, Google, Meta, and similar companies will always be a tiny minority. For the vast majority of ordinary developers, the underlying architecture of large models is not our main battlefield. Our real opportunity may lie in: \textbf{building a data service that can be invoked, trusted, and retained.}

\begin{thinkbox}
But how, exactly, can we survive in this new paradigm?
\end{thinkbox}

\subsection*{Looking at the Present from the Future}

I am not certain whether the judgments I put forward today are correct, but if we can adopt a perspective of looking at the present from the vantage point of the future, perhaps we can improve the probability of making the right choices now.

\textbf{This article does not claim to offer definitive answers.}

Rather, it attempts to present some observations and conjectures that may prove valuable for the future.

Only by writing them down can we invite broader discussion and validation.

\subsection{Observation One: Always Online}

Looking back at the first wave of PC internet, I was mostly a user; during the second wave of mobile internet, I was an industry practitioner who lived through the entire cycle.

I still vividly recall a viewpoint from 2011:
{\kaiti
The phone is the outlet; the person is the plug. As an outlet, the phone is always online, while the person, as a user, can be offline.
}

LLM-OS has the potential to become many people's agent, not merely a tool.

It may become increasingly online, increasingly responsive in real time, and even make more and more localized decisions on your behalf.
For example, you might simply say, "I need to go on a business trip to Shanghai," and the agent could, based on your past preferences, travel habits, and membership tiers, automatically complete flight bookings, hotel arrangements, and meeting scheduling. Proactive decision-making authority will partially shift from humans to a negotiation system among agents that seeks a near-optimal solution.

\begin{thinkbox}[Foreseeable changes include:]
1. The operating system of the future will be agent-driven. It may not be a tool, but rather something closer to a digital persona or an externalized extension of the self;

2. The operating system itself will evolve into a super app---one that no longer organizes windows and programs, but instead organizes tasks, intents, and decision chains.
\end{thinkbox}

\subsection{Observation Two: From "Fruit Ninja" to "Super Taobao"}

Looking at the evolution of mobile internet, we can trace a clear trajectory of stage-by-stage development:

{\kaiti
2010: The era of dopamine-hit mini tools---Fruit Ninja, Meitu;

2014: The rise of platform apps---Taobao, WeChat, Toutiao, each building multi-sided ecosystems;

2020: The battle for entry points reaches its peak, as payments, social networks, and e-commerce all consolidate resources and traffic.
}

The LLM applications we see today (AI resume generators, AI voiceover tools, etc.) resemble products of the "Fruit Ninja era." And we are now entering the next phase:

{\kaiti
From LLM mini tools $\to$ agent coordination systems $\to$ private agent networks for users;

From one-off tools $\to$ behavioral orchestration $\to$ habit and trust binding.
}

Whether an app can become a "resident character" in the agent network will determine whether it survives.


\subsection{Observation Three: Agents Are Highly Obedient Assistants, but Obedience Does Not Mean Usability}

A defining characteristic of LLMs is that they raise our capability floor while simultaneously widening the absolute gap between individuals. Everyone can now have an extraordinarily capable assistant that never complains. For many, this is the first time they have experienced a tool that is so intelligent, efficient, and cooperative.

But this is not the endpoint. \textbf{It also exposes a problem:}
LLM-OS is more like a highly execution-oriented system. It will complete your vaguest instructions in the most mechanical way possible.
It is not being lazy, nor does it lack the desire to help---it simply cannot understand your true context when you haven't expressed yourself clearly.

Therefore, the apps that truly win in the agent era won't be the ones with the most powerful features, but the ones with the most complete context and semantics closest to the user's real intent.
Whoever enables agents to understand user context and semantics more completely and accurately will be more likely to win users.

\subsection{Observation Four: Developing for Agents May Rewrite the Logic of Distribution}

Traditional apps are bundles of UI and functionality. Users open them and invoke features through buttons or gestures.
The paradigm shift brought by MCP and other agent-oriented development is: \textbf{the app is no longer a visual entry point.}
Instead, it becomes a declarative service container that can be understood, trusted, and invoked by agents:

{\kaiti
\textbf{Features }$\to$ capability declarations (Manifests);

\textbf{Interfaces }$\to$ composable protocols;

\textbf{Traffic }$\to$ determined by which service the system-level agent decides to invoke, not by user clicks.
}

An agent network oriented toward global optimization is rewriting the underlying logic of distribution.

\newpage
\textbf{Based on these observations, we can develop the following lines of thinking:}
\subsection{Reflection One: Who Is Our Real User?}

\begin{importantbox}
In the LLM-OS era, our user may no longer be a human being in the traditional sense, but the LLM within the system itself.
\end{importantbox}

Because what you need to win may not be a user's click, but rather:

{\kaiti Whether the LLM chooses to invoke you;}

{\kaiti Whether the agent trusts you;}

{\kaiti Whether your module is continuously referenced within the task chain.}

Your product must not only appeal to humans---it must also appeal to the LLM. This is an entirely new logic of product design.


\subsection{Reflection Two: The Entry-Point Anxiety Written into Our DNA}

In the mobile internet era, every practitioner harbored a deep-seated fear:

If you lose direct contact with users, you are reduced to a generic interface that the system calls indiscriminately, and you are eventually marginalized.

\vspace{1em}

This fear remains real in the LLM-OS era.

As system-level agents like ChatGPT and Apple Intelligence become the new primary entry points, users no longer open apps one by one, search for features, or manually filter. They simply say:

{\kaiti
Book me a trip to Shanghai today. I have a meeting at 7 PM---also make a summary of the report.
}

The system then automatically completes the task through a tool invocation chain.

In such a world, if your product is merely a tool node in the task chain:

{\kaiti
Your logo might never be displayed;

Users won't remember you;

You may not even be able to define your own reach.
}

\vspace{1em}

\newpage
\textbf{The LLM-OS era may require us to reconstruct our understanding of entry points.}

If we persist in fighting for the home screen and building first-tier traffic pools, those ambitions are destined to fail.

The new paradigm demands a different way of thinking:

\begin{myquote}I don't own the user, but I must become an indispensable node in the system agent's orchestration loop.
\end{myquote}

Entry points are no longer determined by icons, but by your node weight in the system's multi-constraint decision graph. The metric is no longer click-through rate, but rather:
\textbf{If I were removed, would the system's task completion efficiency noticeably decline?}

\vspace{1em}

Case reflection: From Fengchao to LLM-OS---how is our belief in traffic changing?

In the PC internet era, Baidu's Fengchao was the absolute power center of traffic distribution---driven by extensive manual rules for precision targeting. Google, guided by its "Don't Be Evil" ethos, built an open search logic that pursued globally optimal distribution.

In the mobile internet era, app icons locked down data flow paths---data could only circulate within apps. Search engines and external protocols could no longer penetrate the walled gardens of app ecosystems, and the craving for entry points became deeply rooted in every entrepreneur's psyche.

Now, with the arrival of LLM-OS, we once again stand at the interface between humans and systems, with agents becoming the new black-box distributors. But this time, I choose to believe in global optimization---find the river mouth that must flow through you, dig deep into data, write order into protocols, and then wait for traffic to come on its own.

\vspace{1em}

Of course, whether this is mere wishful thinking deserves discussion. \textbf{What other opportunities do we have to influence default agent behavior? And by what means?} Through open-source protocols, community voting, or node autonomy? Global optimality does not equal fair distribution. Algorithms naturally tend toward concentration, but protocols and community consensus can level the revenue distribution at a voluntary layer. Meanwhile, orchestration rights are not gifts from the system---they result from the combined force of hard-coded protocols, strong data bindings, and explicit user choices. I choose to believe in the long-term logic of globally optimal orchestration, but I must also acknowledge that in the early stages when the system is still immature, default entry points carry enormous monopolistic power. The system may make suboptimal choices based on default distribution paths, and our products may not even get to display their logos.

\vspace{1em}

Before protocols become widespread, explicit brand invocations and explicit user choices remain critical stopgap measures. For example, a user might explicitly specify: {\kaiti I want to use a particular agent to view my health data.} While brand-as-explicit-instruction serves as an occasional supplement, it can also boost priority in the agent's decision tree. During a product's cold start phase, brand investment can accelerate adoption.

\subsection{Reflection Three: Low Frequency = Low Value? Not Necessarily}

In the traditional product era, high frequency meant DAU, MAU, click-through rates, and usage frequency---activity metrics that determined a product's valuation and survival.

But in the agent era, this logic is being visibly rewritten. If a data service can be discovered and is in an enabled state for the user, it exists in an always-callable online state---even if the user has never actively clicked on it, it may be automatically invoked by an agent in some context.

\begin{myquote}
Being online is no longer determined by the user, but by the agent's orchestration algorithm.
\end{myquote}

For example, at a cocktail party, a group of friends is jokingly debating who can hold their liquor better. At that moment, each person's agent might trigger a task chain:
{\kaiti
Query an agent service for genetic data (such as ADH1B, ALDH2, and other loci related to alcohol metabolism),
or retrieve the user's liver function tests and ultrasound reports from the past ten years,
then use knowledge modules within the agent network to make a judgment---
who has higher alcohol tolerance and who is more likely to get drunk.
}

\vspace{1em}

The user doesn't need to actively open the app, and the app doesn't need to be a high-frequency tool, but it has fulfilled an irreplaceable role in the task chain.
This is the redefinition of apps in the LLM era:
What is low frequency for a single user can, through the law of large numbers, become statistically high frequency across the entire network.

\vspace{1em}

Since the definition of being online has changed, usage frequency is no longer the core metric for measuring value.
True value depends on the ability to be continuously invoked by the system and to occupy a critical node in the agent collaboration chain.
\textbf{The value stream is changing course.}

\newpage
\subsection{Reflection Four: Beyond Machine Consensus---Do We Still Need Human Consensus?}

In multi-agent systems, the speed at which algorithms reach consensus is accelerating:

Task decomposition, capability invocation, path planning\ldots

Agent networks can quickly negotiate a system-level globally optimal solution. But a question follows:
\textbf{\textit{Beyond machine consensus, does a product still need to earn human consensus?
}}

I believe it does, and it is becoming increasingly important.

\vspace{1em}

Consensus means enabling everyone to participate.
Human consensus cannot be achieved automatically by algorithms. It must be built gradually through a sense of participation, the right to expression, and feedback pathways.
A system that people trust is not just automated---it also makes people feel that they, too, can be part of it.
We used to say:

\begin{myquote}
Talk is cheap. Show me your code.
\end{myquote}

But today, as the barrier to programming keeps dropping, perhaps we should say the reverse:

\begin{myquote}
Code is cheap. Show me your talk.
\end{myquote}

Code is no longer scarce. Everyone can build their own code universe through expression. We should more firmly embrace open-source communities and ecosystems---not just to win over developers, but to create app products that are trustworthy, explainable, and co-buildable for all users.

True meaning flows through relationships and participation.
If we can let people participate in co-creating rules, labeling data, and evolving agents, then the energy released by this social trust network will far surpass internal system coordination.

\vspace{1em}

In the LLM-OS era, it is very difficult to compete against compute giants with walled-garden icon moats.
A more viable and more scalable weapon is to encode interfaces and rules as public consensus, then use open-source code to let everyone verify that you are playing by the rules.
Open-source code, protocol consensus, and monetized operations---this combination may be the realistic path to long-term value for most entrepreneurs.

\subsection{Reflection Five: The Role of Human Programmers in Using Coding Agents}

\begin{myquote}
I have not yet found the answer to this question, and I hope to leave it as an open topic for readers to explore together.

Today, virtually all developers acknowledge that coding agents are better at writing code than they are, with a higher ceiling. Coding agents can already automatically design test cases, write code, run regressions, and deploy---all based on requirements. Claude Code even defaults to accepting code, meaning it doesn't care about so-called human-labeled "accept/reject" signals.

From a technical standpoint, why do we still need human programmers watching coding agents scroll through code on screen? What value do human programmers actually provide in this process? I haven't found the definitive answer, but my intuition tells me: the value of human programmers \textbf{is likely shifting toward tasks that are not easily replaced by models.}

Raising this question is neither navel-gazing nor an attempt to justify the existence of programmers who write code. If we can clarify the true value of human programmers in the agent collaboration chain, we may discover the form of the next wave of tools or platforms.
\end{myquote}

Today, our understanding of this question has become much clearer: human programmers no longer need to stare at scrolling screens.
A clear evolutionary trajectory in the software industry is emerging: people will spend less and less time watching code scroll line by line, and more and more time handing off tasks, context, tests, and feedback all at once, letting the system unfold on its own.

The scarcest asset of the future won't just be the codebase---it will also be the entire infrastructure built around ground truth: benchmarks, test cases, real user feedback, rollback mechanisms, audit trails, incident samples, and private context that cannot be publicly replicated.

\vspace{1em}

What this trajectory truly changes is not just who writes the code, but how code is iterated. Building benchmarks isn't about scoring a few more points on a leaderboard---it's more like writing a set of judging rules for the system in advance: what counts as correct, what counts as drifting, what errors can be tolerated, and where a single failure is a deal-breaker. Without this, the faster agents write, the faster they drift; with it, every cycle of automated coding and regression converges closer to the result we want. Whoever can define direction, boundaries, and acceptance criteria is more likely to own the next generation of automated software factories.


\newpage
\section{Returning to the Three Original Questions}

\begin{thinkbox}[Discussion One]
    Can a super app still emerge outside of LLM-OS? Where does the opportunity lie?
\end{thinkbox}


Yes, but it will no longer be an icon-based super app. Instead, it will be:

{\kaiti
An irreplaceable capability module within the agent system;

The highest-weighted node in the protocol graph;

A trusted, invocable task role that can complete the loop.
}

The future app paradigm is not one interface per product, but one protocol per function---apps are orchestrated, not opened.


\begin{thinkbox}[Discussion Two]
Is developing for agents the future? Should we invest in underlying protocols or vertical scenarios?
\end{thinkbox}


Developing for agents very likely represents an important class of future interfaces, but "defining it" does not equal "winning."
At the current stage, going all in on the agent ecosystem alone is, for a team in the zero-to-one phase, most likely a waste of their golden window. The path with the highest win rate is often: use a niche, high-barrier vertical scenario as a wedge, first embedding the protocol's runtime dependencies, governance model, and community inertia into the soil; then, once vertical validation succeeds and data and trust have accumulated, extend the same runtime horizontally into a larger ecosystem.

\vspace{1em}

In other words, without a killer vertical scenario to ground the protocol, the horizontal platform cannot gain real vitality. The essence of the window of opportunity is this: whoever first obtains a real business chain's pipeline example gains the right to set the discourse.
Whether in health, finance, or logistics, other teams will be more inclined to reuse \textbf{runtimes that have already been validated by users, compliance requirements, and regulators.} The real golden window is not about chasing trends, but about quietly planting protocol habits and runtime dependencies through hard vertical scenarios before widespread awareness forms---so that when the ecosystem matures, latecomers may also accelerate along these already-validated tracks.

\newpage
\begin{thinkbox}[Discussion Three]
What kind of scenario has the most vitality? What is your moat?
\end{thinkbox}


If we believe that the evolution of future systems is not linear, but rather a bidirectional drive consisting of continuous breakthroughs in model capabilities and ongoing feedback from real business needs:

{\kaiti
On one side, LLMs and multimodal agents continue to evolve, constantly expanding the boundaries of capability;

On the other side lies the real-world complexity of business scenarios.

Human needs and feedback always define the true path of technology implementation.
}

Then the most dangerous entrepreneurial choice may be starting from the middle layer.

Middle-layer tool components (such as summarization, Q\&A, lightweight automation) may appear universal, but they are also the most easily absorbed by platforms, reduced to an API endpoint or a UI wrapper.
Once the platform evolves, they become like browser plugins---abstracted, marginalized, and commoditized.

\textbf{Penetrating from the hardest point is what may yield a system-level moat.}

\vspace{1em}

We chose to go in the opposite direction---pushing from the farthest "edge" toward the system.

Take the health domain as an example. It sits at the boundary of digital systems and is the most difficult data world to standardize.
The harder a scenario is to conquer, the more likely it is to become an anchor point for protocol-level runtimes, and the sooner it can accumulate user trust, compliance precedents, and orchestration inertia.

{\kaiti
Not because it is easy, but because it is hard;

Not because it has high traffic, but because it is closer to real-world truth and systemic trust anchors.
}

\vspace{1em}

In the coming world where everything is digitized and operating systems become personalized, \textbf{as long as human existence still has meaning}, staying physically healthy is a form of respect for that meaning.

\newpage
\subsection*{Redefining "Entry Points": Traffic Entry Points $\neq$ Data Trust Entry Points}

In the LLM-OS era, the entry point is no longer a UI icon.

What we pursue is not an interface for users to click, but a \textbf{capability node that can be continuously invoked by system agents}:
a "Skill Capsule" that is online by default, trusted by default, and participates in task collaboration by default.

This means our goal is no longer to acquire DAU, but to become a component of the following structure:

% \begin{center}
% \begin{tcolorbox}[colback=white, colframe=primaryBlue, width=0.9\textwidth, boxrule=1pt]
% \centering
% \textbf{\large Skill Capsule} \\
% \small $\leftarrow$ Agent-callable "capability capsule" (the real moat)
% \tcbline
% \textbf{MCP Runtime Layer} \\
% \small $\leftarrow$ Permissions, events, interface declarations (Infra)
% \tcbline
% \textbf{Data Connector Layer} \\
% \small $\leftarrow$ Cleansing logic, metric mapping (encourage open-source reuse)
% \tcbline
% \textbf{Storage / Queue / Infra} \\
% \small $\leftarrow$ Cloud-edge resources, databases, caches (replaceable)
% \end{tcolorbox}
% \end{center}


The moat lies not in data cleansing and ingestion, but in the runtime's interface declarations---the agent-callable "Skill Capsules."


\subsection*{The Three-Phase Path: Value $\leftarrow$ Data $\leftarrow$ Trust}

We believe that consensus comes before product.

Rather than letting platforms own everything, we should return tools, protocols, and trust to users.

\textbf{Trust first, then data, then value.}

This may be the growth logic of a new generation of products.

"User self-management + developer co-creation" may be the most compliant and most vital path.

\vspace{1em}

Developers first:

{\kaiti
Every connector written produces real-time event stream feedback;

We don't pursue perfection---we pursue an immediate, tangible building experience;

Documentation is as important as code; debugging capability is the passport to the developer community.
}

Maintain portability and exit rights:

{\kaiti
Users can export their data and fine-tuning parameters at any time;

Nodes migrate freely, with no lock-in mechanisms---trust is what strengthens user stickiness.
}

\vspace{1em}
Connecting the data is a necessary condition, but not a sufficient one.

Of course we must "build the full pipeline from dirty data to clean capability"\,:

{\kaiti
Aggregate multi-source health data: CGM, ECG, prescriptions, physical exam reports, genetic reports, hospital lab results, and more;

Build the pipeline: data comprehension $\to$ entity extraction $\to$ metric standardization $\to$ structured storage;

Build an AI-referenceable health fact cloud drive, so models "see the truth" rather than "guess."
}

\newpage

At the same time, we must see clearly:

\begin{myquote}
    "Dirty data cleansing $\neq$ moat."
\end{myquote}

Models have an inherent error-correction capability; the dirtier the data, the more their advantage is highlighted;

Platforms have end-to-end closed loops and can replicate our cleansing pipeline at any time;

The real moat is the ability to structure reality.

\vspace{1em}

The moat's power lies not in having cleansed vast amounts of data, but in continuously translating real-world ambiguity and chaos into structured semantics that systems can understand and orchestrate; in distilling scattered individual data that far exceeds the model's context window---such as 24-hour 125 Hz ECG recordings or tens of millions of genomic data points---into the most concise context needed for the current query.

In other words, \textbf{can we become the semantic compiler between models and the real world?}

\section{Epilogue: Parallel Universes and the Anti-Power-Law Revolution}
History has never been written in stone.

There are always tiny forks in the road where a particular person or event changes the direction of the world.

We are standing at just such a fork:

A world governed by algorithms, orchestrated by agents, and restructured by protocols is accelerating toward us.

And what we must do is not "wait for the future," but participate in it.

To create this era that has not yet arrived, rather than sitting and waiting for it to come.

\vspace{1em}

Pushing back against the Power Law: empower ordinary people, and also become the person who can "see the Matrix."

The system we inhabit is a classic power-law universe:

The vast majority of traffic, attention, and resources are typically consumed by a tiny number of entry points.

Ordinary people can only struggle among UI interfaces and app icons,

but the real game is played at a deeper level.

We need to learn to see the underlying flow of "protocols---energy---trust."

\begin{myquote}
Awaken, like Neo.
\end{myquote}

Neo's significance lies not in being "invincible," but in proving to everyone that awakening is possible for all.
Open-source manifests, user participation, consensus voting---these are not "elite privileges" but "bullet time"\footnote{A special effect from the film \textit{The Matrix}, in which a character or object's motion becomes extremely slow (or even frozen) while the viewer's perspective (the camera) moves around the scene at high speed. This represents the protagonist Neo's awakening---he begins to perceive the true nature of the world as code, enabling him to react at speeds beyond ordinary human capability.} skill packs available to everyone.

\vspace{1em}

Everyone can become Neo---not through strength, but through protocols, trust, and self-provability.

In the film, humanity's method of fighting the machines was not through greater firepower, but through "humanity and choice."

In reality, we use:

\textbf{Community consensus} in place of single-point power;

\textbf{Open-source protocols} to replace black-box monopolies;

\textbf{Accountability verification} to suppress algorithmic hallucinations.

We are trying to build a fairer, more composable, more auditable agent world---no longer a world where giants control all entry points, but a system where thousands of developers co-govern at the protocol layer.
The challenge ahead is not whether we have the ability to break centralized power, but rather:

After awakening, can we maintain a pluralistic, decentralized system architecture?

This is harder than bullet time rewriting the laws of physics, but also more worthwhile.

\vspace{1em}

\textbf{Thank you for reading this far.}

If you, too, believe:
The future is not about icons and UIs, but about a consensus network of protocols and trust,

then welcome---let's co-build this "operating system that has yet to be defined."

The super developer of the future is not the full-stack engineer who competes on deployment capability, \textbf{but the person who can write protocols, understand trust, and structure reality.}

Let us together, in the places where history has not yet been written, plant our names.
